\documentclass[11pt]{article}
\usepackage[margin=1in]{geometry}
\usepackage{graphicx}
\usepackage{array}
\usepackage{hyperref}
\usepackage{longtable}
\usepackage{enumitem}
\usepackage{booktabs}
\usepackage{caption}
\usepackage{titlesec}
\titleformat{\section}{\large\bfseries}{\thesection.}{1em}{}

\begin{document}

\begin{center}
    \Large \textbf{Sri Sivasubramaniya Nadar College of Engineering, Chennai} \\
    \large (An autonomous Institution affiliated to Anna University) \\
    \vspace{0.3cm}
\end{center}

\begin{table}[!h]
\renewcommand{\arraystretch}{1.5}
\centering
\resizebox{\textwidth}{!}{%
\begin{tabular}{|l|cll|}
\hline
\textbf{Degree \& Branch}     & \multicolumn{1}{c|}{BE Computer Science \& Engineering} & \multicolumn{1}{l|}{\textbf{Semester}} & VI \\
\hline
\textbf{Subject Code \& Name} & \multicolumn{3}{c|}{ Machine Learning  Laboratory} \\
\hline
\textbf{Academic Year}        & \multicolumn{1}{c|}{2025–2026 } & \textbf{Batch} & 2023–2027 \\
\hline
\end{tabular}
}
\end{table}

\begin{center}
    \textbf{Experiment 8: Clustering Human Activity Recognition Data using K-Means, DBSCAN, and Hierarchical Clustering}
\end{center}

\section*{Objective}
To implement and analyze the performance of clustering algorithms on the Human Activity Recognition dataset:
\begin{itemize}
    \item \textbf{Model A:} K-Means Clustering.
    \item \textbf{Model B:} DBSCAN (Density-Based Spatial Clustering of Applications with Noise).
    \item \textbf{Model C:} Hierarchical Agglomerative Clustering (HAC).
\end{itemize}
Students must visualize and compare the clusters formed by each algorithm, report results, and analyze performance.

\section*{Dataset}
\textbf{Download:} \href{https://archive.ics.uci.edu/dataset/240/human+activity+recognition+using+smartphones}{Human Activity Recognition Using Smartphones Dataset}  

- 30 volunteers (aged 19–48).  
- 6 activities: WALKING, WALKING\_UPSTAIRS, WALKING\_DOWNSTAIRS, SITTING, STANDING, LAYING.  
- Sensors: 3-axis accelerometer and gyroscope (50 Hz).  
- Preprocessed into windows of 2.56s (128 readings/window) with feature extraction from time and frequency domain.  

\section*{Theory}
\subsection*{K-Means Clustering}
- Objective: Partition data into $k$ clusters by minimizing within-cluster variance.  
- Distance metric: Euclidean.  
- Algorithm:
\begin{enumerate}
    \item Initialize $k$ centroids.
    \item Assign points to nearest centroid.
    \item Update centroids as mean of assigned points.
    \item Repeat until convergence.
\end{enumerate}

\subsection*{Elbow Method for Choosing $k$}
- Compute \textbf{Within-Cluster Sum of Squares (WCSS)} for different $k$ values.  
- Plot $k$ vs. WCSS.  
- The ``elbow point'' (where reduction in WCSS slows down) suggests the optimal $k$.  

\begin{table}[!h]
\centering
\caption{K-Means Elbow Method Results}
\renewcommand{\arraystretch}{1.3}
\begin{tabular}{|c|c|c|}
\hline
\textbf{Number of Clusters ($k$)} & \textbf{WCSS (Inertia)} & \textbf{Silhouette Score} \\
\hline
2 &  \_\_\_\_ & \_\_\_\_ \\
3 &  \_\_\_\_ & \_\_\_\_ \\
4 &  \_\_\_\_ & \_\_\_\_ \\
5 &  \_\_\_\_ & \_\_\_\_ \\
6 &  \_\_\_\_ & \_\_\_\_ \\
7 &  \_\_\_\_ & \_\_\_\_ \\
8 &  \_\_\_\_ & \_\_\_\_ \\
\hline
\end{tabular}
\end{table}

\textbf{Students must:}
\begin{itemize}
    \item Fill in the table above after running experiments.
    \item Plot $k$ vs. WCSS (Elbow curve).
    \item Plot $k$ vs. Silhouette Score.
    \item Select and justify the best $k$.
\end{itemize}

\subsection*{DBSCAN}
- Density-based clustering algorithm.  
- Parameters: $\epsilon$ (neighborhood radius), $minPts$ (minimum points).  
- Core points, border points, and noise are identified.  
- Advantage: Can discover clusters of arbitrary shape and handle noise.

\subsection*{Hierarchical Clustering}
- Builds hierarchy of clusters using either:
\begin{itemize}
    \item Agglomerative (bottom-up merging).
    \item Divisive (top-down splitting).
\end{itemize}
- Linkage criteria: Single, Complete, Average, Ward’s method.  
- Represented by a dendrogram.

\section*{Steps for Implementation}
\begin{enumerate}[label=\arabic*.]
    \item Load the dataset and preprocess:
    \begin{itemize}
        \item Handle missing values.
        \item Encode categorical labels.
        \item Apply normalization/standardization.
    \end{itemize}
    \item Perform exploratory data analysis (EDA):
    \begin{itemize}
        \item Visualize feature distributions.
        \item Apply dimensionality reduction (PCA/t-SNE) for visualization.
    \end{itemize}
    \item Apply clustering algorithms:
    \begin{itemize}
        \item \textbf{K-Means:} Use Elbow method to choose $k$.
        \item \textbf{DBSCAN:} Tune $\epsilon$ and $minPts$.
        \item \textbf{Hierarchical:} Apply Agglomerative clustering with Ward’s linkage.
    \end{itemize}
    \item Visualize clusters in 2D using PCA/t-SNE.
    \item Compare results with ground truth activity labels (although clustering is unsupervised).
    \item Analyze performance using clustering metrics.
\end{enumerate}

\section*{Evaluation Metrics}
\begin{itemize}
    \item \textbf{Internal:} Silhouette Score, Davies–Bouldin Index, Calinski–Harabasz Index.
    \item \textbf{External (w.r.t. true labels):} Adjusted Rand Index (ARI), Normalized Mutual Information (NMI).
    \item Confusion matrix (optional, by mapping clusters to activities).
\end{itemize}

\section*{Graphs and Visualization}
Students should produce:
\begin{itemize}
    \item Elbow curve: $k$ vs. WCSS.
    \item Silhouette curve: $k$ vs. Silhouette Score.
    \item Scatter plots of clusters (after PCA/t-SNE).
    \item Dendrograms for hierarchical clustering.
    \item Bar plots of evaluation metrics across algorithms.
\end{itemize}

\section*{Observation Questions}
\begin{itemize}
    \item Which algorithm produced the most meaningful clusters? Why?
    \item How sensitive was K-Means to the choice of $k$?
    \item Did DBSCAN detect noise or small clusters effectively?
    \item How does linkage choice (single/complete/ward) affect hierarchical clustering?
    \item Which internal metric best matched your visual intuition of cluster quality?
\end{itemize}

\section*{Report Checklist}
\begin{itemize}
    \item Aim and Objective
    \item Preprocessing Steps
    \item K-Means with Elbow Method Results
    \item Clustering Implementation (K-Means, DBSCAN, Hierarchical)
    \item Evaluation Metrics and Results
    \item Graphs and Visualizations
    \item Observations and Comparative Analysis
\end{itemize}

\end{document}
